% Luc Destombe --- licence CC BY-NC-SA 4.0
% https://creativecommons.org/licenses/by-nc-sa/4.0/deed.fr
% Fichier autonome : compiler avec pdflatex, deux fois (signets, nombre de pages).
\def\lyceeversion{1.1}
\documentclass[11pt,a4paper]{article}
\makeatletter
\newif\ifeleve
\elevefalse
\elevetrue

\newif\iflycee@palatino
\lycee@palatinotrue

\RequirePackage[utf8]{inputenc}
\RequirePackage[T1]{fontenc}
\RequirePackage[french]{babel}
\RequirePackage{etoolbox}

\newcommand{\ShorthandsOff}{\@ifundefined{shorthandoff}{}{\shorthandoff{;:!?}}}
\newcommand{\ShorthandsOn}{\@ifundefined{shorthandon}{}{\shorthandon{;:!?}}}
\ShorthandsOff
\AtBeginDocument{\ShorthandsOn}
\AtBeginEnvironment{tikzpicture}{\ShorthandsOff}

\RequirePackage{geometry}
\RequirePackage{fancyhdr}
\RequirePackage{lastpage}
\RequirePackage{multicol}
\RequirePackage{array}
\RequirePackage{enumitem}
\RequirePackage{xcolor}
\RequirePackage{needspace}

\RequirePackage{amsmath,amssymb}
\RequirePackage{mathpazo}
\renewcommand{\ttdefault}{lmtt}
\DeclareMathSizes{10.95}{10.95}{8}{5.5}
\begingroup\catcode`\_=\active \catcode`\^=\active
\gdef_#1{\sb{\scriptscriptstyle #1}}
\gdef^#1{\sp{\scriptscriptstyle\lycee@moinsepais #1}}
\endgroup
\newcommand\lycee@moins{\mathbin{%
  \setbox\z@\hbox{$\m@th\scriptscriptstyle\lycee@moinsorig$}%
  \dimen@=.55\wd\z@ \advance\dimen@-.5pt
  \hbox to.75\wd\z@{\hss$\m@th\scriptscriptstyle\vcenter{\hbox{%
    \tikz\draw[line width=.5pt,line cap=round](0,0)--(\the\dimen@,0);}}$\hss}}}
\begingroup\lccode`\~=`\- \lowercase{\endgroup\let~}\lycee@moins
\newcommand\lycee@moinsepais{\mathcode`\-="8000 }
\AtBeginDocument{\mathchardef\lycee@moinsorig=\mathcode`\-\relax}
\AtBeginDocument{\catcode`\_=12 \mathcode`\_="8000
  \catcode`\^=12 \mathcode`\^="8000 }
\newcommand\lycee@exposants{%
  \fontdimen13\textfont2=.48\fontdimen6\textfont2
  \fontdimen14\textfont2=.48\fontdimen6\textfont2
  \fontdimen15\textfont2=.48\fontdimen6\textfont2
  \fontdimen13\scriptfont2=.48\fontdimen6\scriptfont2
  \fontdimen14\scriptfont2=.48\fontdimen6\scriptfont2
  \fontdimen15\scriptfont2=.48\fontdimen6\scriptfont2}
\every@math@size\expandafter{\the\every@math@size\lycee@exposants}
\newlength{\retraitcalculs}
\setlength{\retraitcalculs}{2em}

\RequirePackage{tikz}
\usetikzlibrary{arrows.meta,calc,babel,shapes.misc}
\RequirePackage[most]{tcolorbox}
\tcbuselibrary{skins,breakable}

\definecolor{coulDef}{HTML}{1F4E79}
\definecolor{coulProp}{HTML}{7030A0}
\definecolor{coulRem}{HTML}{C55A11}
\definecolor{coulEx}{HTML}{2E7D32}
\definecolor{coulExo}{HTML}{B03A2E}
\definecolor{coulPart}{HTML}{1F4E79}
\definecolor{coulMeth}{HTML}{1B6E4F}
\definecolor{coulCourbe}{HTML}{0B5ED7}
\definecolor{coulCourbeB}{HTML}{C00000}
\definecolor{coulCourbeC}{HTML}{2E7D32}
\definecolor{coulGrille}{HTML}{8C8C8C}
\definecolor{coulRep}{HTML}{1B6E4F}
\definecolor{coulBar}{HTML}{2E7D32}

\newcommand{\R}{\mathbb{R}}
\newcommand{\Rs}{\mathbb{R}^{*}}
\renewcommand{\le}{\leqslant}
\renewcommand{\ge}{\geqslant}
\renewcommand{\approx}{\simeq}
\newcommand{\vect}[1]{\overrightarrow{#1}}
\newcommand{\norme}[1]{\left\lVert #1 \right\rVert}
\newcommand{\intervalle}[2]{\left[#1\,;#2\right]}
\RequirePackage{cancel}
\renewcommand{\CancelColor}{\color{coulExo}}
\newcommand{\fois}[1]{\times{\color{coulExo}#1}}
\newcommand{\barre}[1]{\cancel{#1}}

\tikzset{grille/.style={coulGrille,line width=0.4pt}}
\newcommand{\quadrillage}[4]{\draw[grille,step=1] (#1,#2) grid (#3,#4);}
\tikzset{
  croix/.style={cross out,draw,line width=0.6pt,inner sep=0pt,minimum size=4pt},
  croixdroite/.style={croix,rotate=45,line width=0.8pt},
}
\newcommand{\pt}[3]{\path (#1) node[croix]{} node[#2,font=\small,inner sep=2.5pt]{$#3$};}

\newcommand{\lycee@repere}[4]{%
  \draw[grille,step=1] (#1,#2) grid (#3,#4);
  \draw[-{Stealth[length=2mm]},thick] (#1,0)--({#3+0.4},0) node[below left=-1pt]{$x$};
  \draw[-{Stealth[length=2mm]},thick] (0,#2)--(0,{#4+0.4}) node[below left=-1pt]{$y$};
  \node[below left,font=\scriptsize,inner sep=1.5pt] at (0,0) {$0$};}
\newcommand{\lycee@gradx}[1]{\foreach \k/\l in {#1}{%
  \draw (\k,0.07)--(\k,-0.07) node[below,font=\scriptsize,inner sep=1.5pt]{$\l$};}}
\newcommand{\lycee@grady}[1]{\foreach \k/\l in {#1}{%
  \draw (0.07,\k)--(-0.07,\k) node[left,font=\scriptsize,inner sep=1.5pt]{$\l$};}}
\newcommand{\lycee@intff}[2]{\left[#1\,;#2\right]}
\newcommand{\lycee@intfo}[2]{\left[#1\,;#2\right[}
\newcommand{\lycee@intof}[2]{\left]#1\,;#2\right]}
\newcommand{\lycee@intoo}[2]{\left]#1\,;#2\right[}
\newcommand{\lycee@ens}[1]{\left\{#1\right\}}
\AtBeginDocument{%
  \providecommand{\repere}{\lycee@repere}%
  \providecommand{\gradx}{\lycee@gradx}%
  \providecommand{\grady}{\lycee@grady}%
  \providecommand{\Cf}{\mathcal{C}\sb{\scriptscriptstyle f}}%
  \providecommand{\Cg}{\mathcal{C}\sb{\scriptscriptstyle g}}%
  \providecommand{\intff}{\lycee@intff}%
  \providecommand{\intfo}{\lycee@intfo}%
  \providecommand{\intof}{\lycee@intof}%
  \providecommand{\intoo}{\lycee@intoo}%
  \providecommand{\ens}{\lycee@ens}}

\newlist{questions}{enumerate}{3}
\setlist[questions,1]{label=\arabic*\textdegree),leftmargin=*,itemsep=2pt,topsep=3pt}
\setlist[questions,2]{label=\alph*),leftmargin=*,itemsep=1pt,topsep=2pt}
\setlist[questions,3]{label=\roman*),leftmargin=*,itemsep=1pt,topsep=1pt}
\setlist[itemize]{label=\textbullet,leftmargin=*,itemsep=2pt,topsep=3pt}

\newcommand{\besoinplace}[1]{\par\penalty\@M\begingroup
  \dimen@=#1\relax\dimen@ii=\pagegoal\advance\dimen@ii-\pagetotal
  \ifdim\dimen@>\dimen@ii\ifdim\dimen@ii>\z@\vfil\fi\break\fi\endgroup}

\newcommand{\lycee@pied}{\thepage/\pageref{LastPage}}
\newcommand{\entete}[2]{%
  \pagestyle{fancy}\fancyhf{}%
  \fancyhead[L]{\footnotesize\color{coulPart}\textbf{#1}}%
  \fancyhead[R]{\footnotesize\color{coulPart}\textbf{#2}}%
  \fancyfoot[C]{\footnotesize\lycee@pied}%
  \renewcommand{\headrulewidth}{0.4pt}%
  \renewcommand{\headrule}{\hbox to\headwidth{%
    \color{coulPart!40}\leaders\hrule height \headrulewidth\hfill}}}

\RequirePackage{ccicons}
\newcommand{\lycee@licence}{{\scriptsize\color{gray}%
  \copyright\ Luc Destombe --- \ccbyncsa\ CC BY-NC-SA 4.0}}
\newif\iflycee@licence \lycee@licencetrue
\newcommand{\sanslicence}{\lycee@licencefalse}
\AtBeginDocument{\iflycee@licence\AddToHookNext{shipout/foreground}{%
  \put(0,0){\rlap{\kern\dimexpr1in+\hoffset+\oddsidemargin+\textwidth\relax
    \llap{\raisebox{-\dimexpr1in+\voffset+\topmargin+\headheight+\headsep
      +\textheight+\footskip\relax}{\lycee@licence}}}}}\fi}

\newcommand{\lycee@site}{faire-des-maths.fr}
\newcommand{\lycee@version}{{\scriptsize\color{gray}\lycee@site
  \ifdefined\lyceeversion\ --- Version~\lyceeversion\fi}}
\AtBeginDocument{\iflycee@licence\AddToHookNext{shipout/foreground}{%
  \put(0,0){\rlap{\kern\dimexpr1in+\hoffset+\oddsidemargin\relax
    \raisebox{-\dimexpr1in+\voffset+\topmargin+\headheight+\headsep
      +\textheight+\footskip\relax}{\lycee@version}}}}\fi}

\newcommand{\formulesserrees}{%
  \setlength{\abovedisplayskip}{3pt}\setlength{\belowdisplayskip}{3pt}%
  \setlength{\abovedisplayshortskip}{2pt}\setlength{\belowdisplayshortskip}{3pt}}

\setlength{\parindent}{0pt}

\RequirePackage[implicit=false,hidelinks,pdfencoding=auto,psdextra,bookmarksopen,
  bookmarksopenlevel=1,pdfcreator={},pdfproducer={}]{hyperref}
\RequirePackage{bookmark}
\hypersetup{pdfauthor={Luc Destombe},pdfkeywords={CC BY-NC-SA 4.0}}
\newcounter{lycee@signet}
\newcommand{\signet}[3][]{\stepcounter{lycee@signet}%
  \edef\lycee@dest{\if\relax\detokenize{#1}\relax
    signet.\arabic{lycee@signet}\else#1\fi}%
  \hypertarget{\lycee@dest}{}\bookmark[level=#2,dest=\lycee@dest]{#3}}
\pdfstringdefDisableCommands{%
  \def\R{R}\def\vect#1{#1}\def\Cf{Cf}\def\Cg{Cg}\def\fois#1{×#1}%
  \def\barre#1{#1}\def\intervalle#1#2{[#1;#2]}\def\intff#1#2{[#1;#2]}%
  \def\textsuperscript#1{#1}\def\dfrac#1#2{#1/#2}\def\frac#1#2{#1/#2}%
  \def\vec#1{#1⃗}}%
\newcommand{\res}[1]{#1}
\geometry{top=2cm,bottom=1cm,left=1cm,right=1cm,headheight=14pt,footskip=0.6cm}

\RequirePackage{pgfplots}
\pgfplotsset{compat=1.18}
\RequirePackage{tkz-tab}

\newcounter{partiecours}
\renewcommand{\thepartiecours}{\Roman{partiecours}}
\newcommand{\partiecours}[1]{%
  \refstepcounter{partiecours}%
  \Needspace*{8\baselineskip}%
  \bigskip
  \noindent\begin{tcolorbox}[enhanced,colback=coulPart,colframe=coulPart,
      boxrule=0pt,arc=2pt,left=8pt,right=8pt,top=4pt,bottom=4pt,
      before skip=6pt,after skip=10pt]
    \color{white}\bfseries\large \thepartiecours{} --- #1%
    \signet[partie-\arabic{partiecours}]{1}{\thepartiecours{} --- #1}
  \end{tcolorbox}%
  \addcontentsline{toc}{section}{\thepartiecours{} --- #1}%
}

\newtcolorbox{boitecours}[3][]{%
  enhanced,breakable,
  colback=white,colframe=#2,coltitle=white,
  fonttitle=\bfseries,
  attach boxed title to top left={xshift=8pt,yshift=-\tcboxedtitleheight/2},
  boxed title style={colback=#2,arc=2pt,boxrule=0pt},
  boxrule=0.9pt,arc=2pt,
  left=8pt,right=8pt,top=8pt,bottom=6pt,
  title={#3},#1}

\newcounter{methode}
\newenvironment{definitionbox}[1][Définition]%
  {\begin{boitecours}[phantom={\signet{2}{#1}}]{coulDef}{#1}}{\end{boitecours}}
\newenvironment{proprietebox}[1][Propriété]%
  {\begin{boitecours}[phantom={\signet{2}{#1}}]{coulProp}{#1}}{\end{boitecours}}
\newenvironment{methodebox}[1][Méthode]{\stepcounter{methode}%
  \begin{boitecours}[phantom={\signet[methode-\arabic{methode}]{2}{#1}}]{coulMeth}{#1}}%
  {\end{boitecours}}

\newcommand{\etape}[1]{\par\smallskip\textbf{\color{coulMeth}#1}\par\nobreak\smallskip}
\newcommand{\enonce}[1]{\emph{#1}\par\smallskip}
\newcommand{\reponse}[1]{\par\smallskip
  {\footnotesize\itshape\color{black!55}Réponses : #1}}
\newcommand{\demo}[1]{\textbf{\color{coulMeth}Démonstration.} #1}
\newcommand{\afaire}[1]{%
  \par\smallskip
  \noindent{\small\color{coulExo}$\blacktriangleright$\ \textbf{À faire :}
    fiche d'exercices du chapitre, #1.}\par\smallskip}

\newenvironment{calculs}{\setlength{\jot}{5pt}\@fleqntrue\@mathmargin\retraitcalculs
  \setlength{\abovedisplayskip}{4pt}\setlength{\belowdisplayskip}{4pt}%
  \setlength{\abovedisplayshortskip}{4pt}\setlength{\belowdisplayshortskip}{4pt}%
  \csname alignat*\endcsname{2}}%
  {\csname endalignat*\endcsname}
\newcommand{\rem}[1]{\text{\small\itshape\color{black!60}#1}}
\newcommand{\explique}[2]{\par\smallskip\noindent
  \begin{minipage}[t]{0.57\linewidth}\raggedright #1\end{minipage}\hfill
  \begin{minipage}[t]{0.40\linewidth}\raggedright\leavevmode
    {\small\itshape\color{black!60}#2}\end{minipage}\par}
\newcommand{\cas}[1]{\par\smallskip\noindent\textbf{\color{coulExo}#1}\vadjust{\nobreak}}
\newcommand{\calc}[1]{\par\smallskip\textbf{\color{coulExo}#1}\par\nobreak}

\newtcolorbox{remarquebox}[1][Remarque]{%
  enhanced,breakable,colback=white,colframe=coulRem,
  boxrule=0pt,leftrule=2.5pt,arc=0pt,outer arc=0pt,
  left=8pt,right=6pt,top=5pt,bottom=5pt,
  before upper={\textbf{\color{coulRem}#1}\par\smallskip}}

\newtcolorbox{exemplebox}[1][Exemple]{%
  enhanced,breakable,colback=white,colframe=coulEx,
  boxrule=0pt,leftrule=2.5pt,arc=0pt,outer arc=0pt,
  left=8pt,right=6pt,top=5pt,bottom=5pt,
  before upper={\textbf{\color{coulEx}#1}\par\smallskip}}

\pgfplotsset{
  repere/.style={
    axis lines=middle,
    axis line style={-{Stealth[length=2.4mm]},thick},
    every axis x label/.style={at={(ticklabel* cs:1.0)},anchor=west},
    every axis y label/.style={at={(ticklabel* cs:1.0)},anchor=south},
    label style={font=\footnotesize},
    tick label style={font=\scriptsize},
    grid=both,
    major grid style={grille},
    minor tick num=0,
    tick align=inside,
    clip mode=individual,
  },
}
\tikzset{
  courbe/.style={coulCourbe,line width=1.1pt,smooth},
  courbeB/.style={coulCourbeB,line width=1.1pt,smooth},
  courbeC/.style={coulCourbeC,line width=1.1pt,smooth},
}

\newlist{etapes}{enumerate}{1}
\setlist[etapes]{label=\textbf{\arabic*.},leftmargin=*,itemsep=1pt,topsep=2pt}

\newlist{expressions}{enumerate}{1}
\setlist[expressions]{label={\Alph*\ $=$},leftmargin=*,itemsep=3pt,topsep=3pt}

\newcounter{exo}
\renewcommand{\theexo}{\ifnum\value{exo}<10 0\fi\arabic{exo}}
\newtcolorbox{exercice}[1][]{%
  enhanced,breakable,
  colback=white,colframe=coulExo,
  boxrule=0.9pt,arc=2pt,
  left=8pt,right=8pt,top=8pt,bottom=6pt,
  coltitle=white,fonttitle=\bfseries,
  attach boxed title to top left={xshift=8pt,yshift=-\tcboxedtitleheight/2},
  boxed title style={colback=coulExo,arc=2pt,boxrule=0pt},
  phantom={\signet{2}{Exercice \arabic{exo}}},
  title={Exercice \theexo\ifx\relax#1\relax\else{} --- #1\fi}}
\newenvironment{exo}[1][\relax]{\refstepcounter{exo}\begin{exercice}[#1]}{\end{exercice}}

  \RequirePackage{comment}
  \excludecomment{correction}
  \renewcommand{\reponse}[1]{}
  \newcommand{\autableau}{\hfill{\footnotesize\itshape\color{black!45}(au tableau)}}
  \newcommand{\etapeexemples}{%
    \par\smallskip\textbf{\color{coulMeth}Exemples}\autableau\par\nobreak\smallskip}
  \newcommand{\etapetableau}[1]{%
    \par\smallskip\textbf{\color{coulMeth}#1}\autableau\par\nobreak\smallskip}
  \newtcolorbox{exempletableau}[1][Exemple]{%
    enhanced,breakable,colback=white,colframe=coulEx,
    boxrule=0pt,leftrule=2.5pt,arc=0pt,outer arc=0pt,
    left=8pt,right=6pt,top=4pt,bottom=4pt,
    before skip=5pt,after skip=5pt,
    before upper={\textbf{\color{coulEx}#1}\autableau\par\smallskip}}

\RequirePackage{comment}
\excludecomment{correctionavj}

\newcommand{\coursEntete}{}
\newcommand{\coursNiveau}{}
\newcommand{\coursTitre}{}
\newcommand{\coursSurTitre}{}
\newcommand{\coursSousTitre}{}

\newcommand{\enteteCours}[1]{\renewcommand{\coursEntete}{#1}}
\newcommand{\chapitrecours}[2]{\enteteCours{Chapitre #1 --- #2}}
\newcommand{\niveaucours}[1]{\renewcommand{\coursNiveau}{#1}}
\newcommand{\titrecours}[3][]{%
  \renewcommand{\coursTitre}{#2}%
  \renewcommand{\coursSurTitre}{#1}%
  \renewcommand{\coursSousTitre}{#3}}

\pagestyle{fancy}
\fancyhf{}
\fancyhead[L]{\footnotesize\color{coulPart}\textbf{\coursEntete}}
\fancyhead[R]{\footnotesize\color{coulPart}\textbf{\coursNiveau}}
\fancyfoot[C]{\footnotesize\thepage}
\renewcommand{\headrulewidth}{0.4pt}
\renewcommand{\headrule}{\hbox to\headwidth{\color{coulPart!40}\leaders\hrule height \headrulewidth\hfill}}

\setlength{\parskip}{2pt}

\newcommand{\lycee@mentioneleve}{\\[2pt]{\small\itshape Version élève}}
\AtBeginDocument{%
  \begin{center}
    \begin{tcolorbox}[enhanced,width=0.72\linewidth,colback=white,colframe=coulPart,
        boxrule=1.6pt,arc=3pt,halign=center,top=10pt,bottom=10pt]
      \ifx\coursSurTitre\empty
        {\Huge\bfseries\color{coulPart} \coursTitre}\\[4pt]
      \else
        {\Huge\bfseries\color{coulPart} \coursTitre}\\[3pt]
        {\Large\bfseries\color{coulPart} \coursSurTitre}\\[5pt]
      \fi
      {\normalsize \coursSousTitre}%
      \lycee@mentioneleve
    \end{tcolorbox}
  \end{center}%
}

\RequirePackage{tikz}
\RequirePackage{xcolor}

\usetikzlibrary{decorations.text}
\definecolor{lycee@metal}{HTML}{1F4E79}
\definecolor{lycee@manche}{HTML}{8A5A2B}
\definecolor{lycee@etincelle}{HTML}{F28C28}
\newcommand{\lycee@pluriel}[2]{#1 #2\ifnum#1>1 s\fi}
\newcommand{\lycee@sceau}{%
  \edef\lycee@bas{\lycee@pluriel{\lyceerelectures}{relecture} et
    \lycee@pluriel{\lyceerodages}{rodage}}%
  \begin{tikzpicture}[line join=round]
    \fill[coulPart!10,draw=coulPart,line width=.8pt] (0:2.2)
      \foreach \i in {1,...,71}{ -- (\i*5:{isodd(\i)?2.1:2.2}) } -- cycle;
    \draw[coulPart,line width=1.2pt] (0,0) circle (1.93);
    \draw[coulPart,line width=.4pt] (0,0) circle (1.85);
    \draw[coulPart,line width=.4pt] (0,0) circle (1.2);
    \path[decorate,decoration={text along path,text align=center,
        text={|\fontsize{9.5}{11}\selectfont\bfseries\color{coulPart}|TEST{\'E} AVEC MES {\'E}L{\`E}VES}}]
      (180:1.45) arc (180:0:1.45);
    \path[decorate,decoration={text along path,text align=center,
        text/.expanded={|\noexpand\fontsize{8.5}{10}\noexpand\selectfont\noexpand\color{coulPart}|\lycee@bas}}]
      (195:1.68) arc (195:345:1.68);
    \begin{scope}[x=.25mm,y=-.25mm]
      \begin{scope}[shift={(-32,-40)}]%
        \fill[lycee@metal] (23,31) rectangle (41,49);
        \draw[lycee@etincelle,line width=.5pt,line cap=round]
          (20,30)--(15,27) (19,34)--(14,35) (44,30)--(49,27) (45,34)--(50,35);
        \fill[lycee@manche,rounded corners=.3mm] (38,24.6) rectangle (60,27.6);
        \fill[lycee@metal,rounded corners=.4mm] (25,20) rectangle (39,31);
        \fill[lycee@etincelle,rounded corners=.2mm] (25,20) rectangle (28,31);
      \end{scope}
    \end{scope}
    \ifdefined\lyceesceauplusieurs\def\lycee@annee{Ann{\'e}es~:}\else
      \def\lycee@annee{Ann{\'e}e}\fi
    \path[decorate,decoration={text along path,text align=center,
        text/.expanded={|\noexpand\fontsize{8}{10}\noexpand\selectfont\noexpand\color{coulPart}|%
          \unexpanded\expandafter{\lycee@annee} \unexpanded\expandafter{\lyceesceauannees}}}]
      (95:1.08) arc (95:445:1.08);
  \end{tikzpicture}}
\AtEndDocument{\ifdefined\lyceesceauannees
  \par\besoinplace{5cm}\vspace{1.5em}%
  \begin{center}\lycee@sceau\end{center}\fi}
\makeatother

\chapitrecours{10}{Fonction exponentielle}
\niveaucours{1\textsuperscript{re} spécialité mathématiques}
\titrecours{FONCTION EXPONENTIELLE}{Cours --- Première spécialité}

\begin{document}

\newcommand{\e}{\mathrm{e}}

\begin{remarquebox}[Comment utiliser ce cours]
Chaque partie commence par ce qu'il faut \textbf{savoir} (définitions, propriétés), puis
vient ce qu'il faut \textbf{savoir faire} : une méthode, avec des
\textbf{exemples} faits en classe, puis des exercices
\og\textbf{À vous de jouer}\fg{} à chercher seul.

\smallskip
Sur cette feuille, les exemples n'ont que leur \textbf{énoncé} : la résolution, marquée
\emph{(au tableau)}, s'écrit dans le cahier pendant le cours. Les réponses des « À vous
de jouer » ne sont pas données.

\end{remarquebox}

\partiecours{Définition, relation fonctionnelle, notation $\e^{x}$}

\begin{remarquebox}[D'où vient cette fonction ?]
Dans une culture de bactéries, plus il y a de bactéries, plus elles se multiplient vite :
la vitesse de croissance est proportionnelle au nombre de bactéries. Avec une fonction
$f$, cela s'écrit $f'=f$ (dans les bonnes unités) : la dérivée est égale à la fonction
elle-même.
\end{remarquebox}

\begin{definitionbox}[Définition 1 --- Fonction exponentielle]
Il existe une seule fonction $f$ dérivable sur $\R$ telle que
\[
  f'=f \qquad\text{et}\qquad f(0)=1 \qquad\text{(admis).}
\]
C'est la \textbf{fonction exponentielle}, notée $\exp$. Ainsi $\exp(0)=1$ et, pour tout
nombre réel $x$, $\exp'(x)=\exp(x)$.
\end{definitionbox}

\begin{proprietebox}[Propriété 1 --- Relation fonctionnelle]
Pour tous nombres réels $x$ et $y$ :
\[
  \exp(x+y)=\exp(x)\times\exp(y) \qquad\text{(admis)}\qquad\text{et}\qquad
  \exp(x)\times\exp(-x)=1.
\]
En particulier, $\exp(x)$ n'est jamais nul.
\end{proprietebox}

\begin{remarquebox}[Démonstration de $\exp(x)\times\exp(-x)=1$]
On pose $\varphi(x)=\exp(x)\times\exp(-x)$. La dérivée de $\exp(-x)$ est
$-\exp(-x)$ (composée, avec $u(x)=-x$ et $u'(x)=-1$).
\begin{calculs}
\varphi'(x) &= \exp(x)\times\exp(-x)+\exp(x)\times\big(-\exp(-x)\big) &\qquad& \rem{$(uv)'=u'v+uv'$ et $\exp'=\exp$}\\
\varphi'(x) &= 0 &&
\end{calculs}
La fonction $\varphi$ est donc constante (ch. 4) ; elle vaut
$\varphi(0)=\exp(0)\times\exp(0)=1$.
\end{remarquebox}

\begin{definitionbox}[Définition 2 --- Le nombre e et la notation $\e^{x}$]
Le nombre $\e=\exp(1)$ est appelé \textbf{nombre d'Euler} ; à la calculatrice,
$\e\approx 2{,}718$.

Pour tout nombre réel $x$, on note $\exp(x)=\e^{x}$ (touche $\e^{x}$ de la
calculatrice).
\end{definitionbox}

La notation est cohérente avec les puissances : $\exp(2)=\exp(1+1)=\e\times\e$, et
$\e\times\e=\e^{2}$. Les règles de calcul sont celles des puissances de $10$, vues au
collège ($10^{3}\times 10^{2}=10^{5}$).

\begin{proprietebox}[Propriété 2 --- Règles de calcul]
Pour tous nombres réels $a$ et $b$, et tout entier relatif $n$ :
\begin{center}
\renewcommand{\arraystretch}{2.4}\setlength{\tabcolsep}{14pt}
\begin{tabular}{|c|c|c|c|}
\hline
$\e^{a}\times\e^{b}=\e^{a+b}$ & $\dfrac{\e^{a}}{\e^{b}}=\e^{a-b}$ &
$\dfrac{1}{\e^{a}}=\e^{-a}$ & $\left(\e^{a}\right)^{n}=\e^{na}$\\ \hline
\end{tabular}
\end{center}
De plus, $\e^{0}=1$ et $\e^{1}=\e$.
\end{proprietebox}

\begin{proprietebox}[Propriété 3 --- Signe]
Pour tout nombre réel $x$, $\e^{x}>0$.
\end{proprietebox}

\begin{remarquebox}[Démonstration]
\begin{calculs}
\e^{x} &= \e^{\frac{x}{2}+\frac{x}{2}} &&\\
\e^{x} &= \left(\e^{\frac{x}{2}}\right)^{2} &\qquad& \rem{$\e^{a}\times\e^{a}=\left(\e^{a}\right)^{2}$}
\end{calculs}
Un carré est positif ou nul, et $\e^{x}$ n'est jamais nul (propriété 1) : $\e^{x}>0$.
\end{remarquebox}

\begin{methodebox}[Méthode 1 --- Simplifier une écriture avec des exponentielles]
\etapeexemples
Pour tout nombre réel $x$, écrire chaque expression sous la forme $\e^{\dots}$, ou
développer.
\begin{questions}[label=\alph*)]
  \item $A=\e^{4}\times\e^{-1}$
  \item $B=\dfrac{\e^{5}}{\e^{-2}}$
  \item $C=\left(\e^{-3}\right)^{2}\times\e^{8}$
  \item $D=\e^{3x}\times\e^{1-x}$
  \item $E=\dfrac{\e^{2x+1}}{\e^{x-3}}$
  \item $F=\left(\e^{x}+2\right)^{2}$ (développer)
\end{questions}

\etape{À vous de jouer}
Pour tout nombre réel $x$, écrire chaque expression sous la forme $\e^{\dots}$, ou
développer.
\begin{center}
$A=\e^{6}\times\e^{-4}$ \qquad $B=\dfrac{\e^{3}}{\e^{-5}}$ \qquad
$C=\left(\e^{-2}\right)^{3}\times\e^{7}$\\[4pt]
$D=\e^{5x}\times\e^{2-3x}$ \qquad $E=\dfrac{\e^{3x-2}}{\e^{x+1}}$ \qquad
$F=\left(\e^{x}+3\right)^{2}$
\end{center}

\end{methodebox}

\afaire{exercices 1 à 5}

\partiecours{Équations et inéquations avec des exponentielles}

\begin{proprietebox}[Propriété 4 --- Comparer deux exponentielles]
La fonction exponentielle est \textbf{strictement croissante} sur $\R$. Par conséquent,
pour tous nombres réels $a$ et $b$ :
\[
  \e^{a}=\e^{b} \iff a=b \qquad\qquad \e^{a}<\e^{b} \iff a<b.
\]
\end{proprietebox}

\begin{remarquebox}[Pourquoi ?]
La fonction exponentielle est strictement croissante (sa dérivée, $\exp$ elle-même, est
strictement positive d'après la propriété 3), donc elle conserve l'ordre :
$4<5$ et donc $\e^{4}<\e^{5}$ ; à l'inverse, $\e^{a}<\e^{b}$ seulement si $a<b$.

Si deux nombres ont la même image par $\exp$, ces deux nombres sont égaux :
$\e^{a}=\e^{b}$ seulement si $a=b$.

Pour comparer à $1$, on écrit $1=\e^{0}$. Une équation $\e^{a}=-5$ n'a pas de solution,
puisqu'une exponentielle est toujours strictement positive.
\end{remarquebox}

\begin{methodebox}[Méthode 2 --- Résoudre une équation avec des exponentielles]
\etapeexemples
Résoudre dans $\R$ chaque équation.
\begin{questions}[label=\alph*)]
  \item $\e^{2x-1}=\e^{x+3}$
  \item $\e^{3x-6}=1$
  \item $\e^{x^{2}}=\e^{x+6}$
  \item $\e^{x}\times\e^{x+4}=\e^{10}$
  \item $\e^{2x}=-5$
\end{questions}

\etape{À vous de jouer}
Résoudre dans $\R$ chaque équation.
\begin{center}
$\e^{3x+2}=\e^{x+8}$ \qquad $\e^{4x+8}=1$ \qquad $\e^{x^{2}}=\e^{2x+3}$\\[4pt]
$\e^{2x}\times\e^{x-1}=\e^{8}$ \qquad $\e^{x+1}=-3$
\end{center}

\end{methodebox}

\begin{methodebox}[Méthode 3 --- Résoudre une inéquation avec des exponentielles]
\etapeexemples
Résoudre dans $\R$ chaque inéquation.
\begin{questions}[label=\alph*)]
  \item $\e^{3x-2}\le\e^{x+4}$
  \item $\e^{5-2x}>1$
  \item $\e^{x^{2}}<\e^{4x-3}$
\end{questions}

\etape{À vous de jouer}
Résoudre dans $\R$ chaque inéquation.
\begin{center}
$\e^{4x-1}\le\e^{x+5}$ \qquad $\e^{6-3x}>1$ \qquad $\e^{x^{2}}<\e^{5x-4}$
\end{center}

\end{methodebox}

\afaire{exercices 6 à 10}

\partiecours{Étude de la fonction exponentielle}

\begin{minipage}[c]{0.58\linewidth}\raggedright
\begin{proprietebox}[Propriété 5 --- Variations et courbe]
La fonction exponentielle est dérivable sur $\R$, et
${\left(\e^{x}\right)'=\e^{x}}$.

Elle est strictement positive et strictement croissante sur $\R$ :
\begin{center}
\begin{tikzpicture}[scale=0.85]
  \tkzTabInit[lgt=1.6,espcl=4.4]{$x$/0.8,$\left(\e^{x}\right)'$/0.8,$\e^{x}$/1.6}{$-\infty$,$+\infty$}
  \tkzTabLine{,+,}
  \tkzTabVar{-/{},+/{}}
  \tkzTabVal{1}{2}{0.4}{$0$}{$1$}
  \tkzTabVal{1}{2}{0.7}{$1$}{$\e$}
\end{tikzpicture}
\end{center}
\end{proprietebox}
\end{minipage}\hfill
\begin{minipage}[c]{0.38\linewidth}\centering
\begin{tikzpicture}[x=0.75cm,y=0.75cm]
  \repere{-3}{-1}{2}{4}
  \gradx{-2,-1,1} \grady{1,2,3}
  \draw[coulCourbe,line width=1.1pt,domain=-3:1.38,samples=80]
    plot (\x,{exp(\x)});
  \draw[coulExo,line width=0.8pt,domain=-2:1.6] plot (\x,{\x+1});
  \draw[coulCourbeC,line width=0.8pt,domain=0.2:1.45] plot (\x,{2.71828*\x});
  \path[coulCourbe] (0,1) node[croix]{} (1,2.71828) node[croix]{};
  \node[coulCourbe,font=\footnotesize] at (-2.2,0.55) {$y=\e^{x}$};
  \node[coulExo,font=\footnotesize] at (1.9,2.2) {$T_{0}$};
  \node[coulCourbeC,font=\footnotesize] at (1.75,3.6) {$T_{1}$};
\end{tikzpicture}
\end{minipage}

\medskip
La courbe passe par les points de coordonnées $(0\,;1)$ et $(1\,;\e)$ ; elle est
au-dessus de l'axe des abscisses, et s'en rapproche vers la gauche. Ses tangentes en
ces deux points, de pente $\e^{0}=1$ et $\e^{1}=\e$, ont pour équations
$T_{0}$ : $y=x+1$ et $T_{1}$ : $y=\e x$. Vers la droite, la fonction croît de plus en
plus vite : $\e^{10}$ dépasse $22\,000$.

\begin{remarquebox}[Rappel du chapitre 4]
Les formules $(u+v)'=u'+v'$, $(ku)'=ku'$, $(uv)'=u'v+uv'$ et
$\left(\dfrac{u}{v}\right)'=\dfrac{u'v-uv'}{v^{2}}$ s'appliquent avec
$\left(\e^{x}\right)'=\e^{x}$. Dans un produit, on factorise ensuite par $\e^{x}$.
\end{remarquebox}

\begin{methodebox}[Méthode 4 --- Dériver une fonction qui contient $\e^{x}$]
\etapeexemples
Calculer la dérivée de chaque fonction ; factoriser par $\e^{x}$ quand c'est possible.
\begin{questions}[label=\alph*)]
  \item $f(x)=3\e^{x}-x^{2}$, sur $\R$
  \item $g(x)=(2x+1)\e^{x}$, sur $\R$
  \item $h(x)=\dfrac{\e^{x}}{x+1}$, sur $\R\setminus\ens{-1}$
\end{questions}

\etape{À vous de jouer}
Calculer la dérivée de chaque fonction ; factoriser par $\e^{x}$ quand c'est possible.
\begin{center}
$f(x)=5\e^{x}-x^{3}$ \qquad $g(x)=(3x-2)\e^{x}$ \qquad
$h(x)=\dfrac{\e^{x}}{x-2}$ sur $\R\setminus\ens{2}$
\end{center}

\end{methodebox}

\afaire{exercices 11 à 14}

\partiecours{Dérivée de $\e^{u}$}

Au chapitre 4, on a dérivé une composée : si $f(x)=v\big(u(x)\big)$, alors
$f'(x)=u'(x)\times v'\big(u(x)\big)$. Avec $v=\exp$, qui est sa propre dérivée, on
garde l'exponentielle telle quelle et on multiplie par $u'(x)$.

\begin{proprietebox}[Propriété 6 --- Dérivée de $\e^{u}$]
Si $u$ est dérivable sur un intervalle $I$, la fonction $x\mapsto\e^{u(x)}$ est
dérivable sur $I$ et
\[
  \left(\e^{u}\right)'=u'\,\e^{u}.
\]
En particulier, pour tous nombres réels $a$ et $b$ :
$\left(\e^{ax+b}\right)'=a\,\e^{ax+b}$ ; pour une variable $t$,
$\left(\e^{at}\right)'=a\,\e^{at}$.
\end{proprietebox}

\begin{methodebox}[Méthode 5 --- Dériver une fonction de la forme $\e^{u}$]
\etapeexemples
Calculer la dérivée de chaque fonction, définie sur $\R$.
\begin{questions}[label=\alph*)]
  \item $f(x)=\e^{4x}$
  \item $g(x)=\e^{-3x+2}$
  \item $h(t)=200\,\e^{-0{,}05t}$
  \item $k(x)=\e^{x^{2}-2x}$
  \item $m(x)=x\,\e^{-2x}$
\end{questions}

\etape{À vous de jouer}
Calculer la dérivée de chaque fonction, définie sur $\R$.
\begin{center}
$f(x)=\e^{5x}$ \qquad $g(x)=\e^{-2x+7}$ \qquad $h(t)=300\,\e^{-0{,}02t}$\\[4pt]
$k(x)=\e^{x^{2}+4x}$ \qquad $m(x)=x\,\e^{-3x}$
\end{center}

\end{methodebox}

\afaire{exercices 15 à 18}

\partiecours{Étude de fonctions avec des exponentielles}

\begin{remarquebox}[Le signe d'une dérivée avec une exponentielle]
Une exponentielle est toujours strictement positive (propriété 3). Si
$f'(x)=(ax+b)\,\e^{x}$, alors $f'(x)$ a le \textbf{signe de $ax+b$} : on étudie le signe
du seul facteur qui peut changer de signe. Pour le reste, le plan d'étude d'une fonction
est celui du chapitre 4.
\end{remarquebox}

\begin{methodebox}[Méthode 6 --- Étudier une fonction qui contient une exponentielle]
\etapeexemples
On considère la fonction $f$ définie sur $\intff{-3}{2}$ par $f(x)=(x-1)\e^{x}$.
\begin{questions}[label=\alph*)]
  \item Calculer $f'(x)$.
  \item Étudier le signe de $f'(x)$.
  \item Dresser le tableau de variations de $f$ et donner son minimum.
  \item Donner une équation de la tangente $T$ à $\Cf$ au point d'abscisse $1$.
  \item Compléter un tableau de valeurs (au dixième, pas de $1$), puis tracer $T$ et
        $\Cf$.
\end{questions}

\etape{À vous de jouer}
On considère la fonction $g$ définie sur $\intff{-2}{3}$ par $g(x)=(x-3)\e^{x}$. Mêmes
questions a) à e), avec la tangente au point d'abscisse $1$.

\end{methodebox}

\afaire{exercices 19 à 22}

\partiecours{Croissance et décroissance exponentielles}

\begin{minipage}[c]{0.58\linewidth}\raggedright
\begin{proprietebox}[Propriété 7 --- Fonctions $t\mapsto\e^{kt}$ et $t\mapsto\e^{-kt}$]
Soit $k>0$.
\begin{itemize}
  \item $\left(\e^{kt}\right)'=k\,\e^{kt}>0$ : $t\mapsto\e^{kt}$ est strictement
        croissante ; c'est une \textbf{croissance exponentielle}.
  \item $\left(\e^{-kt}\right)'=-k\,\e^{-kt}<0$ : $t\mapsto\e^{-kt}$ est strictement
        décroissante ; c'est une \textbf{décroissance exponentielle}.
\end{itemize}
Les deux courbes passent par le point de coordonnées $(0\,;1)$. Plus $k$ est grand,
plus la variation est rapide.
\end{proprietebox}
\end{minipage}\hfill
\begin{minipage}[c]{0.38\linewidth}\centering
\begin{tikzpicture}
\begin{axis}[repere,
    x=0.6cm,y=0.6cm,
    xmin=-3.4,xmax=3.6,ymin=-0.6,ymax=4.6,
    xtick={-3,-2,-1,1,2,3},ytick={1,2,3,4},
    xlabel={$t$},ylabel={$y$}]
  \addplot[courbe,domain=-3:1.38,samples=60]{exp(x)};
  \addplot[courbeB,domain=-3:2.8,samples=60]{exp(0.5*x)};
  \addplot[courbeC,domain=-1.38:3,samples=60]{exp(-x)};
  \node[coulCourbe,font=\footnotesize,left] at (axis cs:1.3,4.2) {$\e^{t}$};
  \node[coulCourbeB,font=\footnotesize,right] at (axis cs:2.4,2.4) {$\e^{0{,}5t}$};
  \node[coulCourbeC,font=\footnotesize,left] at (axis cs:-1.4,4.2) {$\e^{-t}$};
\end{axis}
\end{tikzpicture}
\end{minipage}

\medskip
Au chapitre 7, une évolution à taux fixe est modélisée par une suite géométrique :
$u_{n+1}=q\times u_{n}$, où $q=1+t$ est le coefficient multiplicateur (une baisse de
$8\,\%$ donne $q=1-0{,}08$, soit $q=0{,}92$). L'exponentielle prolonge ce modèle aux
instants qui ne sont pas entiers.

\begin{proprietebox}[Propriété 8 --- Lien avec les suites géométriques]
Soit $C$ et $k$ deux nombres réels, et $u_{n}=C\,\e^{kn}$ pour tout entier naturel $n$.
Alors $(u_{n})$ est une \textbf{suite géométrique} de premier terme $u_{0}=C$ et de raison
$q=\e^{k}$.
Si $k>0$, alors $q>1$ ; si $k<0$, alors $0<q<1$.
\end{proprietebox}

\begin{remarquebox}[Démonstration]
\begin{calculs}
u_{n+1} &= C\,\e^{k(n+1)} &\qquad& \rem{on remplace $n$ par $n+1$}\\
u_{n+1} &= C\,\e^{kn+k} &\qquad& \rem{on développe l'exposant}\\
u_{n+1} &= C\,\e^{kn}\times\e^{k} &\qquad& \rem{$\e^{a+b}=\e^{a}\times\e^{b}$}\\
u_{n+1} &= u_{n}\times\e^{k} &\qquad& \rem{$C\,\e^{kn}=u_{n}$}
\end{calculs}
\end{remarquebox}

\begin{methodebox}[Méthode 7 --- Tracer les courbes de $t\mapsto\e^{kt}$ et $t\mapsto\e^{-kt}$]
\etapeexemples
On considère les fonctions $f$ et $g$ définies sur $\intff{-3}{3}$ par
$f(t)=\e^{0{,}5t}$ et $g(t)=\e^{-0{,}5t}$.
\begin{questions}[label=\alph*)]
  \item Donner le sens de variation de $f$ et celui de $g$.
  \item Compléter un tableau de valeurs de $f$ et de $g$ (au dixième, pas de $1$).
  \item Tracer les deux courbes dans un même repère.
\end{questions}

\etape{À vous de jouer}
Mêmes questions pour $f(t)=\e^{0{,}4t}$ et $g(t)=\e^{-0{,}4t}$, sur $\intff{-3}{3}$.

\end{methodebox}

\begin{methodebox}[Méthode 8 --- Passer d'un modèle exponentiel à une suite géométrique]
\etapeexemples
Un laboratoire reçoit un échantillon de $500$~mg d'iode radioactif. La masse restante,
en mg, $t$ jours plus tard, est $m(t)=500\,\e^{-0{,}086t}$.
\begin{questions}[label=\alph*)]
  \item Calculer $m(0)$ et interpréter le résultat.
  \item Donner le sens de variation de $m$.
  \item On note $u_{n}=m(n)$ la masse après $n$ jours. Montrer que $(u_{n})$ est une suite
        géométrique ; donner sa raison, puis son arrondi au millième.
  \item De quel pourcentage la masse baisse-t-elle chaque jour ?
  \item Calculer la masse restante après $8$ jours, arrondie au mg.
\end{questions}

\etape{À vous de jouer}
Une culture contient $2\,000$ bactéries. Leur nombre $t$ heures plus tard est
$N(t)=2\,000\,\e^{0{,}3t}$.
\begin{questions}
  \item Calculer $N(0)$ et interpréter le résultat.
  \item Donner le sens de variation de $N$.
  \item On note $u_{n}=N(n)$ le nombre de bactéries après $n$ heures. Montrer que
        $(u_{n})$ est une suite géométrique ; donner sa raison, puis son arrondi au
        millième.
  \item De quel pourcentage le nombre de bactéries augmente-t-il chaque heure ?
  \item Calculer le nombre de bactéries après $4$ heures, arrondi à l'unité.
\end{questions}

\end{methodebox}

\afaire{exercices 23 à 26}

\partiecours{Problèmes concrets}

\begin{remarquebox}[Dans un sujet d'épreuve]
L'énoncé donne souvent la dérivée (« montrer que $f'(t)=\dots$ ») : on calcule $f'(t)$
et on arrive exactement à l'expression écrite. Une valeur approchée se lit à la
calculatrice (touche $\e^{x}$) ; une durée cherchée en heures entières se lit dans un
tableau de valeurs, en s'appuyant sur le sens de variation. On conclut toujours par une
phrase, avec l'unité.
\end{remarquebox}

\begin{methodebox}[Méthode 9 --- Étudier une fonction exponentielle dans un contexte]
\etapeexemples
Une vidéo est mise en ligne à midi. Le nombre de personnes qui la regardent en même
temps, en milliers, est donné $t$ heures plus tard par
\[
  V(t)=10t\,\e^{-0{,}5t},\qquad t\in\intff{0}{10}.
\]
\begin{questions}[label=\alph*)]
  \item Calculer $V(0)$ et interpréter le résultat.
  \item Montrer que $V'(t)=(10-5t)\,\e^{-0{,}5t}$.
  \item Étudier le signe de $V'(t)$, puis dresser le tableau de variations de $V$.
  \item À quelle heure le nombre de spectateurs est-il maximal ? Combien sont-ils, à la
        centaine près ?
  \item Voici des valeurs de $V(t)$ arrondies au centième.
        \begin{center}
        \renewcommand{\arraystretch}{1.3}
        \begin{tabular}{|c|c|c|c|c|c|}
        \hline
        $t$ & $5$ & $6$ & $7$ & $8$ & $9$\\ \hline
        $V(t)$ & $4{,}10$ & $2{,}99$ & $2{,}11$ & $1{,}47$ & $1{,}00$\\ \hline
        \end{tabular}
        \end{center}
        À partir de quelle heure entière y a-t-il moins de $2\,000$ spectateurs ?
\end{questions}

\etape{À vous de jouer}
Un patient prend un médicament à 8~h. La concentration du médicament dans le sang, en
mg/L, est donnée $t$ heures plus tard par
\[
  C(t)=8t\,\e^{-0{,}25t},\qquad t\in\intff{0}{12}.
\]
\begin{questions}
  \item Calculer $C(0)$ et interpréter le résultat.
  \item Montrer que $C'(t)=(8-2t)\,\e^{-0{,}25t}$.
  \item Étudier le signe de $C'(t)$, puis dresser le tableau de variations de $C$.
  \item À quelle heure la concentration est-elle maximale ? Que vaut-elle, au centième
        près ?
  \item Voici des valeurs de $C(t)$ arrondies au centième.
        \begin{center}
        \renewcommand{\arraystretch}{1.3}
        \begin{tabular}{|c|c|c|c|c|c|}
        \hline
        $t$ & $8$ & $9$ & $10$ & $11$ & $12$\\ \hline
        $C(t)$ & $8{,}66$ & $7{,}59$ & $6{,}57$ & $5{,}63$ & $4{,}78$\\ \hline
        \end{tabular}
        \end{center}
        Le médicament agit tant que la concentration est d'au moins $6$~mg/L. À partir
        de quelle heure entière n'agit-il plus ?
\end{questions}

\end{methodebox}

\afaire{exercices 27 à 29}

\end{document}
